%=====================================================================
% Cost-Sensitive Fraud Decisioning Under Delayed and Censored Labels
% IEEE Conference Template (IEEEtran)
%=====================================================================
%
% HOW TO USE THIS FILE ON OVERLEAF
% --------------------------------
% 1. Create a new Overleaf project (blank project).
% 2. Replace the default main.tex with this file.
% 3. Compile (pdfLaTeX). IEEEtran is preinstalled on Overleaf.
%
% PLACEHOLDERS TO FILL BEFORE SUBMISSION
% --------------------------------------
% [A] Author email: search "email@students.national-u.edu.ph"
% [B] Repository URL in Appendix B: search "REPOSITORY-URL-HERE"
% [C] Verify final page count is within 6–10 (excluding appendices).
% [D] VERIFY Streamlit URL in Appendix B opens without error.
%
%=====================================================================

\documentclass[conference]{IEEEtran}

\usepackage{cite}
\usepackage{amsmath,amssymb,amsfonts}
\usepackage{booktabs}
\usepackage{graphicx}
\usepackage{textcomp}
\usepackage{xcolor}
\usepackage{url}
\usepackage[hidelinks]{hyperref}

% Roman section numbering is the IEEEtran default; bold top-level headings (instructor's format)
\usepackage{etoolbox}
\makeatletter
\patchcmd{\section}{\scshape}{\scshape\bfseries}{}{\typeout{WARNING: section-heading bold patch failed}}
\patchcmd{\subsection}{\itshape}{\itshape\bfseries}{}{\typeout{WARNING: subsection-heading bold patch failed}}
\makeatother

% Correct bad hyphenation here
\hyphenation{op-tical net-works semi-conduc-tor}

\begin{document}

%---------------------------------------------------------------------
% TITLE AND AUTHORS
%---------------------------------------------------------------------
\title{Cost-Sensitive Fraud Decisioning Under Delayed and Censored Labels: A Temporal Backtest on the Bank Account Fraud Dataset}

\author{
\IEEEauthorblockN{Ken Ira Lacson Talingting}
\IEEEauthorblockA{College of Computing and Information Technologies\\
National University Philippines, Manila, Philippines\\
talingtingkl@students.national-u.edu.ph}
\and
\IEEEauthorblockN{Lawrence Abad}
\IEEEauthorblockA{College of Computing and Information Technologies\\
National University Philippines, Manila, Philippines\\
abadlp@students.national-u.edu.ph}
\and
\IEEEauthorblockN{John Kelly Nogales}
\IEEEauthorblockA{College of Computing and Information Technologies\\
National University Philippines, Manila, Philippines\\
nogalesjd@students.national-u.edu.ph}
\and
\IEEEauthorblockN{Smith Albert Surio}
\IEEEauthorblockA{College of Computing and Information Technologies\\
National University Philippines, Manila, Philippines\\
suriosb@students.national-u.edu.ph}
}

\maketitle

%---------------------------------------------------------------------
% ABSTRACT AND INDEX TERMS (run-in style, per sample)
%---------------------------------------------------------------------
\begin{abstract}
Fraud detection systems must decide whether to approve, review, or block a
transaction in real time, yet fraud labels arrive weeks to months later and
are frequently censored before the evaluation window closes. This paper
frames the problem as a \textbf{cost-sensitive decision problem} rather than
a classification problem and evaluates it under a \textbf{temporal backtest}
on the Bank Account Fraud (BAF) dataset with a simulated one-month label
delay. A Logistic Regression classifier feeding an
\textbf{argmin-of-expected-cost policy} reduces realized cost per
transaction by \textbf{57.91\%} relative to the strongest baseline,
LightGBM + static 0.5 at \textbf{0.017798}, with a 95\% bootstrap confidence
interval of \textbf{[53.95\%, 62.13\%]}. The result is robust to a
2$\times$ variation in each cost parameter and does not depend on test-set
tuning: the cost rate is derived from training-window amounts only.
Censored-label counts (96,843 of 1,000,000, or 9.68\%) are reported rather
than treating unobserved fraud as legitimate, and the full evaluation
protocol is released alongside the pipeline. The paper contributes a
reproducible template for evaluating fraud decision policies under delayed
and partial labels, and demonstrates that static thresholds on imbalanced
data behave near-identically to approve-all even when the underlying model
is well-calibrated. Limitations include a single delay regime, a synthetic
dataset, and an amount proxy derived from proposed credit limit.
\end{abstract}

\begin{IEEEkeywords}
Cost-sensitive learning, fraud detection, delayed feedback, censored
labels, decision policy, temporal backtest, Logistic Regression, Bank
Account Fraud dataset, expected-cost minimization, false-positive cost,
bootstrap confidence intervals, sensitivity analysis, machine learning.
\end{IEEEkeywords}

%=====================================================================
\section{Introduction}
%=====================================================================

Digital payment systems in the Philippines process millions of
transactions every day. According to the Bangko Sentral ng Pilipinas
(BSP), the share of digital payments in total retail payment volume rose
from 20.1\% in 2020 to 52.8\% in 2023, surpassing the central bank's
50\% target~\cite{bsp2023roadmap}. The value of transactions processed through Philippine payment systems
has grown substantially over the same period~\cite{bspspei}. This
growth has brought fraud into closer view: the BSP's Consumer Assistance
Mechanism processed over 10,000 consumer complaints in 2023, with
unauthorized online transfers and fraudulent withdrawals representing a
substantial share~\cite{bspcam}. For banks, issuers, merchants, and their
customers, the operational question is not whether a transaction is
fraudulent in hindsight but whether to approve, review, or block it
\emph{before} the label arrives.

The fraud label itself is the core difficulty. When a customer's card is
used without authorization, the loss is not confirmed at the moment of the
transaction---it is confirmed weeks to months later, when the customer
files a dispute, the merchant responds, and the chargeback is resolved.
Industry dispute windows commonly run 30 to 120 days~\cite{visa}. On
real-time payment rails such as InstaPay and PESONet, funds may become
irrevocable before fraud is confirmed. A machine-learning model trained on
settled disputes therefore learns from a label that arrives late---and
from a label that may never arrive at all, because some fraud goes
unreported, some disputes are resolved in the merchant's favor, and some
victims do not notice~\cite{chapelle2014}.

Most published work in fraud detection treats the problem as binary
classification and reports ranking metrics such as area under the ROC
curve (AUC) or area under the precision-recall curve (PR-AUC). These
metrics are useful for comparing models on a fixed labeled dataset, but
they do not answer the operational question a bank actually faces: given
the model's score, should this transaction be approved, reviewed, or
blocked, and what is the realized cost of that decision? On a dataset
where fraud is roughly 1\% of transactions, a high AUC can coexist with a
policy that barely improves on approving everything.

This paper takes a different framing. Fraud detection is treated as a
one-step decision problem under delayed, partial, and asymmetric-cost
feedback. The model produces a fraud probability, but the system under
evaluation is the \emph{policy} that maps that probability, the transaction
amount, and a fixed cost matrix into one of three actions. The policy
chooses the action that minimizes expected cost. This is not a novel
policy---expected-cost minimization is standard in cost-sensitive
learning~\cite{elkan2001}---but it is rarely evaluated end-to-end on a
fraud dataset with simulated label delay, censored labels reported rather
than hidden, and statistical rigor on the primary metric.

The framing is evaluated on the Bank Account Fraud (BAF) dataset, a
public 1-million-row synthetic benchmark designed for fraud
research~\cite{jesus2022}. Because BAF provides only month-level time
granularity, a one-month label delay is simulated and the final month is
held out as censored. A cost-sensitive policy is compared against six
baselines under an identical temporal split and cost matrix. Realized cost
per transaction is the primary metric. Robustness is tested with a
2$\times$ sensitivity analysis on each cost parameter and a bootstrap
confidence interval.

The main findings are as follows. First, the cost-sensitive policy reduces
realized cost per transaction by \textbf{57.91\%} relative to the strongest
baseline (LightGBM + static 0.5), with a 95\% bootstrap confidence interval
of \textbf{[53.95\%, 62.13\%]}. Second, the static threshold performs
\emph{near-identically to approve-all} (5.97\% reduction), a failure mode
that AUC-based evaluation would not reveal. Third, the result is robust
across a 2$\times$ range on each of three cost parameters; the minimum
advantage across all variations is \textbf{47.76\%}. Fourth, the cost rate
used for amount scaling is derived from training-window amounts only, so
the headline result does not depend on any test-set tuning.

This paper makes three contributions. It contributes a reproducible
evaluation protocol for fraud decision policies under delayed and censored
labels, including pre-registered stop criteria and an explicit prohibition
on random splits. It contributes a temporal backtest on BAF with six
canonical baselines, censored-label reporting, cost sensitivity, and
bootstrap confidence intervals. It also contributes a documented negative
result at the framing level: static thresholds on highly imbalanced fraud
data are not decision rules, and reporting them without a cost anchor
overstates model performance.

The remainder of this paper is organized as follows. Section~II reviews
related work in cost-sensitive learning, delayed feedback, and fraud
detection. Section~III describes the decision policy and cost formulation.
Section~IV describes the BAF dataset, the delay simulation, and the
temporal split. Section~V describes the experimental setup and baselines.
Section~VI presents results, sensitivity analysis, and confidence
intervals. Section~VII discusses what the results imply for real fraud
operations. Section~VIII states the limitations. Section~IX concludes and
outlines future work.

%=====================================================================
\section{Related Work}
%=====================================================================

Fraud decisioning sits at the intersection of three research traditions:
cost-sensitive learning, learning under delayed or partial feedback, and
imbalanced fraud detection. Each contributes a piece of the framing adopted
here, and each leaves a gap that this paper addresses.

\subsection{Cost-Sensitive Learning}

Cost-sensitive learning reframes classification around the asymmetric cost
of errors rather than the symmetric error rate~\cite{elkan2001}. Elkan
formalized the Bayes-optimal decision rule under an arbitrary cost matrix,
showing that the optimal threshold on the posterior probability is a
function of the cost ratio between false positives and false negatives, not
a fixed value like 0.5~\cite{elkan2001}. The framework was later extended
to settings where only positive labels are observed and the unlabeled
population must be modeled explicitly~\cite{elkan2008}. In fraud detection
the cost asymmetry is severe: approving a fraudulent transaction exposes
the issuer to the full fraud loss, while blocking a legitimate one incurs
customer friction and churn~\cite{elkan2001,dalpozzolo2015}. A classifier
trained to optimize accuracy or F1 treats these errors as symmetric and
therefore selects the wrong operating point. This is the foundation of the
decision-centric framing adopted here: the classifier is an input to a
policy, and the policy chooses the action that minimizes expected cost.

\subsection{Learning Under Delayed and Partial Feedback}

Chapelle modeled delayed feedback in display advertising as a survival
problem, using survival analysis to estimate the conversion rate under an
unknown delay distribution rather than assuming the label is either present
or absent~\cite{chapelle2014}. The fraud setting shares the structural
property---a decision must be made before the outcome is known---but
differs in that fraud labels are frequently censored rather than merely
delayed. When the observation window closes before a chargeback resolves,
the label is missing, not negative~\cite{jesus2022}. Treating censored
labels as negatives biases the classifier toward under-detection and
inflates apparent precision. The final month of the dataset is therefore
held out as censored, excluded from training and evaluation, and the
censored count is reported alongside every result.

\subsection{Fraud Detection and Calibration}

Dal Pozzolo et al.\ showed that probability calibration degrades under
severe class imbalance and that uncalibrated scores undermine any
downstream cost-sensitive threshold~\cite{dalpozzolo2015}. Their result
matters directly for the present work: the decision policy evaluated here
consumes a calibrated fraud probability, and the expected cost of each
action is linear in that probability. An uncalibrated score produces wrong
expected costs, and therefore wrong decisions. Calibration is therefore
treated as a gate---ECE below 0.05 on the validation window---rather than a
diagnostic. A classifier that fails the gate is rejected as a policy input
regardless of its ranking performance.

\subsection{Benchmark Datasets for Fraud and Imbalanced Learning}

The Bank Account Fraud (BAF) dataset was introduced by Jesus et al.\ as a
realistic, temporal, imbalanced benchmark for machine learning
evaluation~\cite{jesus2022}. It includes a month index for chronological
splitting, a mixed-type feature set, and documented biases that make it
suitable for stress-testing evaluation protocols. This work uses BAF's v1
dataset because its month-level granularity supports a defensible delay
simulation, its fraud rate of approximately 1.1\% forces the cost asymmetry
to matter, and its licence (CC~BY~4.0) and citable source make the results
reproducible.

\subsection{The Gap This Paper Fills}

Existing work tends to evaluate either classifiers or policies, but rarely
both, and even more rarely under a simulated delay with censored labels
explicitly reported. Studies on imbalanced fraud data often optimize
ranking metrics without connecting them to an operational cost; studies on
cost-sensitive learning often assume the label is available at decision
time; studies on delayed feedback often focus on advertising, where the
delay distribution is well understood, rather than on fraud, where labels
are both late and frequently missing. This paper bridges these threads by
evaluating an argmin-of-expected-cost policy end-to-end on BAF under a
simulated one-month label delay, with censored labels counted rather than
hidden, with a pre-registered evaluation protocol, and with bootstrap
confidence intervals on the primary metric. The classifier is treated as an
input to the policy; the policy is the object of study.

%=====================================================================
\section{Methodology}
%=====================================================================

This section describes the decision policy that is the object of study,
the cost matrix it consumes, the three classifiers that are compared as
policy inputs, the preprocessing applied to each, the hyperparameter search
procedure, and the calibration gate that a classifier must pass before it
can feed the policy.

\subsection{Decision Formulation}

At decision time, the system must choose one of three actions for each
transaction: \texttt{approve}, \texttt{review}, or \texttt{block}. Three
actions are used because approve and block alone ignore manual review,
which is where most operational fraud decisioning lives, and because a
three-action formulation exposes the review-cost tradeoff directly. For a
transaction with predicted fraud probability $p = P(\text{fraud} \mid X_t)$,
the policy computes the expected cost of each action and selects the
minimum:

\begin{equation}
\begin{aligned}
\mathbb{E}[C_{\text{app}}] &= p \cdot \ell(a), \\
\mathbb{E}[C_{\text{rev}}] &= c_r + p \cdot \ell_r, \\
\mathbb{E}[C_{\text{blk}}] &= (1 - p) \cdot c_f,
\end{aligned}
\end{equation}

where $a$ is the transaction amount, $\ell(a)$ is the amount-scaled fraud
loss, $c_r$ is the review cost, $\ell_r$ is the residual fraud loss after
review, and $c_f$ is the false-positive cost. The policy selects:

\begin{equation}
a^{*} = \arg\min_{a \in \{\text{app},\,\text{rev},\,\text{blk}\}} \mathbb{E}[C_a].
\end{equation}

The argmin rule is the source of truth. Two thresholds on $p$ can be
derived from the cost matrix as a diagnostic view, but they are not the
implementation. Under amount scaling, the review threshold becomes per-row,
and the fixed-threshold view no longer describes the policy's behavior. The
argmin is always well-defined, including in the degenerate cases where the
thresholds collapse or fall outside $[0,1]$; these edge cases are covered
by unit tests.

\subsection{Cost Matrix and Amount Scaling}

All cost inputs are frozen before any modeling begins and are never tuned
on validation or test data. The values are $\text{fraud\_loss} = 1.0$,
$\text{false\_positive\_cost} = 0.1$, $\text{review\_cost} = 0.02$, and
$\text{residual\_fraud\_loss} = 0.3$. The constant-loss matrix treats fraud
loss as a fixed quantity of 1.0. This is a simplification: a fraudulent
transaction of 10,000 pesos costs more than one of 10 pesos. The policy
therefore supports amount-scaled fraud loss, in which the fraud-loss term
becomes per-row:

\begin{equation}
\ell(a) = a \cdot r_{\ell},
\end{equation}

where the rate is calibrated from training-window amounts only:

\begin{equation}
r_{\ell} = \frac{1}{\overline{a}_{\text{train}}}
         = \frac{1}{521.1626}
         = 0.0019187869.
\end{equation}

This keeps the mean of $r_{\ell} \cdot a$ equal to 1.0 on the training
window, matching the constant-loss comparison scale, without using any
test-window information. The amount column used is
\texttt{amount\_proxy}, which is derived from \texttt{proposed\_credit\_limit}
in the Bank Account Fraud dataset; a true transaction-amount column is not
available, and this is documented as a limitation.

\subsection{The Three Classifiers}

Three traditional classifiers are compared as policy inputs.
\textbf{Logistic Regression} is included as a linear, interpretable
baseline~\cite{pedregosa2011}. It is configured with $C = 10.0$ and
\texttt{max\_iter}$=1000$. \textbf{Random Forest} is included as a
non-linear ensemble that handles mixed feature types and interactions
without explicit feature engineering, configured with
\texttt{n\_estimators}$=400$, \texttt{n\_jobs}$=-1$, and
\texttt{random\_state}$=42$. \textbf{LightGBM} is included as a
gradient-boosting method with strong tabular performance and native
categorical handling~\cite{ke2017}, configured with
\texttt{n\_estimators}$=300$, \texttt{learning\_rate}$=0.05$, and
\texttt{random\_state}$=42$. No class weighting is applied to any
classifier. The decision policy assumes $p$ is a calibrated probability,
and reweighting to correct for class imbalance would distort that
probability. Cost asymmetry is handled by the policy, not by the training
objective.

\subsection{Preprocessing Design}

Preprocessing is fit within each training fold and applied unchanged to the
corresponding validation and test data. Logistic Regression uses
StandardScaler on numeric features and one-hot encoding on the five
categorical features with \texttt{handle\_unknown="ignore"}. Random Forest
uses no scaling and ordinal encoding with
\texttt{handle\_unknown="use\_encoded\_value"} and
\texttt{unknown\_value}$=-1$. LightGBM uses native categorical handling via
pandas \texttt{category} dtype, with the category mapping captured during
training and reapplied at inference. Eight columns are excluded from all
classifiers: \texttt{transaction\_id}, \texttt{fraud\_bool},
\texttt{month}, \texttt{label\_month}, \texttt{observed},
\texttt{amount\_proxy}, \texttt{proposed\_credit\_limit}, and
\texttt{device\_fraud\_count}. The result is a 28-feature input: 23 numeric
and 5 categorical.

\subsection{Hyperparameter Search}

Hyperparameter selection is performed by expanding-window cross-validation
on the training window. Each fold uses contiguous months for training and
the next month for validation; no shuffling occurs and no fold mixes time
boundaries. On the current training window (months 0--2), this yields two
folds. Preprocessing is refit within each fold's training portion. Logistic
Regression tests four values of $C$ in $\{0.01, 0.1, 1.0, 10.0\}$. Random
Forest tests three forest sizes in $\{100, 200, 400\}$. LightGBM tests four
hand-selected combinations: $\{100, 0.10\}$, $\{300, 0.10\}$,
$\{300, 0.05\}$, and $\{500, 0.05\}$, where each pair is
(\texttt{n\_estimators}, \texttt{learning\_rate}). The selected
configuration per classifier is the one with the lowest mean realized cost
after the policy across folds. The selection is subject to a pre-registered
\textbf{noise-band guard}: the top classifier must win every fold and the
relative gap to the runner-up must exceed 5\%. If either condition fails,
the comparison is reported as a non-finding and the tie is broken in favor
of the simplest classifier under the ordering Logistic Regression $>$
Random Forest $>$ LightGBM.

\subsection{Calibration Gate}

The decision policy assumes $p$ is a probability: expected costs are linear
in $p$, and a score that is not a probability produces wrong expected costs
and therefore wrong actions~\cite{dalpozzolo2015}. A \textbf{calibration
gate} is therefore enforced before any classifier can be admitted as a
policy input. The gate requires an Expected Calibration Error (ECE) below
0.05 on the validation window, computed over 10 quantile bins. If a
classifier passes the gate, no calibration step is applied. If it fails,
exactly one calibration method is fit on validation and re-measured.
Calibration is never fit on the test window.

\subsection{Evaluation Criterion}

The classifier is selected by realized cost after the policy on the
validation window, following the same cost matrix and the same delay regime
as the final test evaluation. The selected classifier is retrained on the
full training window and evaluated once on the untouched test window
(months 5--6). The test window is used exactly once. All model selection,
threshold diagnostics, and calibration checks happen before that
evaluation; nothing is tuned on the test window.

%=====================================================================
\section{Dataset and Preprocessing}
%=====================================================================

\subsection{Dataset}

The Bank Account Fraud (BAF) Suite, version 1, is a public 1-million-row
synthetic benchmark introduced by Jesus et al.\ for machine learning
evaluation under realistic fraud conditions~\cite{jesus2022}. The dataset
has 32 raw columns including the binary target \texttt{fraud\_bool}, a
month index ranging from 0 to 7, and a mix of numeric, categorical, and
binary features representing an account-application transaction. The
overall fraud rate is 1.1029\%. The dataset is distributed under
CC~BY~4.0.

\subsection{Delay Simulation}

BAF exposes only month-level time granularity. Under a one-month delay
regime, a transaction in month $t$ has its label available at month
$t+1$. A transaction in month 7 has \texttt{label\_time} $=8$, which is
beyond the observation window; it is \emph{censored}, not negative. Under
this regime, 96,843 rows (9.68\% of the dataset) are censored and are
excluded from training and evaluation. Their count is reported in every
results table as a data-integrity figure.

\subsection{Chronological Split}

The dataset is split chronologically into training months 0--2 (397,039
rows, fraud rate 0.9813\%), validation months 3--4 (278,627 rows, fraud
rate 1.0207\%), and test months 5--6 (227,491 rows, fraud rate 1.2576\%).
No shuffling occurs. No fold mixes time boundaries. Preprocessing is fit on
the training window only. The test window is used exactly once.

\subsection{Preprocessing Implementation}

Preprocessing is applied per classifier and fit within each training fold.
The categorical columns are \texttt{payment\_type},
\texttt{employment\_status}, \texttt{housing\_status}, \texttt{source}, and
\texttt{device\_os}. Numeric features are used as-is for tree-based
classifiers and standardized for Logistic Regression. The full
preprocessing pipeline is saved alongside the selected classifier so that
the deployed application applies the exact same transformations used
during training.

%=====================================================================
\section{Experimental Setup}
%=====================================================================

\subsection{Evaluation Protocol}

The primary metric is realized cost per transaction on the test window
under the frozen cost matrix and the one-month delay regime. The policy is
compared against the strongest canonical baseline, not the weakest.
Statistical claims use a 95\% bootstrap confidence interval on the
advantage, with 1,000 resamples and seed 42. The minimum effect size for a
finding is a 5\% relative reduction.

\subsection{Baselines}

Six canonical baselines are reported: random decision, approve-all,
block-all, and each of the three classifiers under a static 0.5 threshold
(Logistic Regression + 0.5, Random Forest + 0.5, LightGBM + 0.5). Every
baseline uses the same split, the same cost matrix, and the same delay
regime. No baseline is tuned on the test window.

\subsection{Sensitivity Analysis}

Each cost parameter (\texttt{false\_positive\_cost}, \texttt{review\_cost},
\texttt{residual\_fraud\_loss}) is varied one at a time over a 2$\times$
range, and the backtest is recomputed at each setting. The model is not
retrained; only the cost matrix changes. The minimum advantage across the
sweep is reported as the robustness bound. The fourth cost parameter,
\texttt{fraud\_loss}, is tested via the constant-vs-amount-scaled
comparison in Section~III-B rather than a 2$\times$ range, because its role
is structural: it determines whether amount scaling is active, not a
numeric magnitude.

%=====================================================================
\section{Results}
%=====================================================================

This section reports the data and split actually used, the cross-validated
classifier comparison, the cost-sensitive policy on the untouched test
window, the statistical and sensitivity analyses that establish whether the
result is a finding or a non-finding, and the operational action
distribution produced by the policy.

\subsection{Data and Split}

The dataset contains 1,000,000 transactions over eight monthly periods
with an overall fraud rate of 1.1029\%. Under the one-month delay regime,
the final month is censored, yielding 96,843 censored rows, or 9.68\% of
the dataset. The single chronological split assigns months 0--2 to
training, months 3--4 to validation, and months 5--6 to test. The test
fraud rate (1.2576\%) is modestly higher than the training rate (0.9813\%),
a known temporal drift documented as a limitation.

Exploratory analysis surfaced three findings that shaped preprocessing and
evaluation. First, the fraud rate is non-monotonic across months: it begins
at 1.13\% in month 0, dips to 0.87\% by month 2, then rises steadily to
1.47\% by month 7---a temporal drift pattern consistent with the BAF
dataset documentation~\cite{jesus2022}. Second,
\texttt{proposed\_credit\_limit} is heavily right-skewed (skew 1.30, mean
515.85, median 200) with 24.17\% of values above the IQR upper bound,
justifying its exclusion from classification and its use as a cost-scaling
input instead. Third, the strongest point-biserial correlations with the
target are \texttt{credit\_risk\_score} ($r=0.0706$) and
\texttt{proposed\_credit\_limit} ($r=0.0689$), both weak but consistent
with fraud detection being a many-weak-signals problem.

\subsection{Classifier Comparison}

Three classifiers were trained under identical folds, preprocessing logic,
and evaluation criterion. Hyperparameters were selected per algorithm by
mean realized cost after the policy across two expanding-window folds.

\begin{table}[!t]
\caption{Cross-validated realized cost per transaction after the policy.}
\label{tab:cv}
\centering
\begin{tabular}{lr}
\toprule
Classifier & Mean $\pm$ SD (cost/txn) \\
\midrule
LightGBM & $0.005852 \pm 0.000300$ \\
Logistic Regression & $0.005915 \pm 0.000338$ \\
Random Forest & $0.006515 \pm 0.000579$ \\
\bottomrule
\end{tabular}
\end{table}

LightGBM achieves the lowest mean cost, with Logistic Regression second and
Random Forest clearly worse. The gap between LightGBM and Logistic
Regression is 0.000063 in absolute terms, or 1.07\% relative---well
inside the 5\% noise-band guard. The comparison is therefore reported as a
\textbf{non-finding}, and the selection is broken by simplicity under the
ordering Logistic Regression $>$ Random Forest $>$ LightGBM. Logistic
Regression is selected as the policy input and retrained on the full
training window.

\begin{table}[!t]
\caption{Logistic Regression classification behavior at the 0.5 cut on the
test window (supporting evidence only).}
\label{tab:cls}
\centering
\begin{tabular}{lr}
\toprule
Metric & Value \\
\midrule
Macro F1 & 0.5031 \\
Accuracy & 0.9875 \\
Precision (fraud) & 0.7826 \\
Recall (fraud) & 0.0063 \\
ROC-AUC & 0.8753 \\
Confusion matrix (TN, FP, FN, TP) & (224,625, 5, 2,843, 18) \\
\bottomrule
\end{tabular}
\end{table}

The low recall at 0.5 reflects the base rate: at a 1.26\% fraud rate, a
threshold that minimizes misclassification produces very few positive
predictions. This is precisely the failure mode the decision-centric
framing is designed to avoid. The 0.5 cut is not the deployed operating
point.

\subsection{Decision Policy on the Test Window}

The selected classifier's calibrated probability feeds the
argmin-of-expected-cost policy under the frozen cost matrix and the
training-calibrated amount-scaling rate. The policy's realized cost per
transaction is \textbf{0.007491}. The strongest baseline is \textbf{LightGBM
+ static 0.5} at \textbf{0.017798}. The policy reduces realized cost per
transaction by \textbf{57.91\%} relative to that baseline.

\begin{table}[!t]
\caption{Policy vs.\ baseline realized cost per transaction (test window).}
\label{tab:baselines}
\centering
\begin{tabular}{lr}
\toprule
Baseline & Cost/txn \\
\midrule
Random decision & 0.047484 \\
Approve-all & 0.018928 \\
Block-all & 0.098742 \\
Logistic Regression + static 0.5 & 0.018737 \\
Random Forest + static 0.5 & 0.018928 \\
LightGBM + static 0.5 (strongest) & 0.017798 \\
\textbf{Cost-sensitive policy (this work)} & \textbf{0.007491} \\
\bottomrule
\end{tabular}
\end{table}

Two observations matter for interpretation. First, the strongest baseline
is LightGBM + static 0.5, not Logistic Regression + static 0.5. The
advantage is reported against the stronger baseline. Second, the three
static-0.5 baselines perform near-identically to approve-all: LightGBM +
0.5 reduces cost by 5.97\%, Logistic Regression + 0.5 by 1.01\%, and
Random Forest + 0.5 by 0\%. A well-calibrated classifier with a static
threshold is effectively not a decision rule at this base rate.

The calibration gate was passed on the validation window with an ECE of
\textbf{0.0033}, well below the 0.05 threshold. The test-window Brier score
for the policy's probability input is \textbf{0.011506}.

\subsection{Statistical Validity}

The advantage was tested with 1,000 bootstrap resamples of the test window
(seed 42) to produce a 95\% confidence interval on the difference between
the baseline and the policy. The point estimate of the difference is
0.010307 per transaction, with a 95\% CI of [0.009602, 0.011058]. The
corresponding advantage is \textbf{57.91\%}, with a 95\% CI of
\textbf{[53.95\%, 62.13\%]}. The interval excludes zero, so the
pre-registered success criterion is met.

\subsection{Sensitivity to Cost Assumptions}

\begin{table*}[!t]
\caption{Cost-parameter sensitivity of the policy advantage.}
\label{tab:sens}
\centering
\begin{tabular}{lrrr}
\toprule
Configuration & Policy cost/txn & Strongest baseline & Advantage \\
\midrule
Baseline & 0.007491 & LGBM + 0.5 (0.017798) & 57.91\% \\
\texttt{false\_positive\_cost} = 0.05 & 0.007037 & LGBM + 0.5 (0.017752) & 60.36\% \\
\texttt{false\_positive\_cost} = 0.20 & 0.007625 & LGBM + 0.5 (0.017890) & 57.38\% \\
\texttt{review\_cost} = 0.01 & 0.006253 & LGBM + 0.5 (0.017798) & 64.87\% \\
\texttt{review\_cost} = 0.04 & 0.009298 & LGBM + 0.5 (0.017798) & \textbf{47.76\%} \\
\texttt{residual\_fraud\_loss} = 0.15 & 0.006588 & LGBM + 0.5 (0.017798) & 62.99\% \\
\texttt{residual\_fraud\_loss} = 0.60 & 0.008700 & LGBM + 0.5 (0.017798) & 51.12\% \\
\bottomrule
\end{tabular}
\end{table*}

The minimum advantage across the sweep is \textbf{47.76\%}, attained at
\texttt{review\_cost} $=0.04$. Even at that worst case the policy retains
nearly half the baseline cost reduction.

\subsection{Operational Action Distribution}

On the 227,491 test transactions, the policy routes 205,395 to
\texttt{approve}, 21,371 to \texttt{review}, and 725 to \texttt{block}.
The operational precision and recall of the policy are 0.0654 and 0.5051
respectively. Ranking metrics at fixed budgets are Precision@1\%~$=0.2431$,
Precision@5\%~$=0.1189$, Precision@10\%~$=0.0775$, and
Recall@1\%~$=0.1933$, Recall@5\%~$=0.4729$, Recall@10\%~$=0.6162$. The
ranking metrics are modest in absolute terms, which is expected under a
1.26\% base rate. The result of interest is the cost reduction, not the
ranking.

%=====================================================================
\section{Discussion}
%=====================================================================

\subsection{Why the Classifier Comparison Is a Non-Finding}

The classifier comparison was deliberately structured so that a
cross-validated ranking could only be reported as a finding if it was
robust to noise. LightGBM achieved the lowest mean cross-validated realized
cost against Logistic Regression by 1.07\%, which is well inside the range
that a two-fold cross-validation on a 1.1\% positive-rate problem can
produce by chance. The comparison is therefore reported as a non-finding,
and the selection is broken in favor of the simpler, more interpretable
classifier. Reporting a 1.07\% difference as substantive would overclaim
on evidence that does not support the claim.

The practical consequence is that a fraud operations team considering a
deployment on this data would not need to prefer a gradient-boosting model
over logistic regression on performance grounds alone. The
interpretability of a linear model---coefficients are readable, decisions
are traceable, drift is diagnosable---carries real operational value that
a 1.07\% cost difference does not outweigh.

\subsection{Why the Policy Beats the Static Baselines}

The 57.91\% cost reduction comes from the decision rule, not from the
classifier. All three classifiers were well-calibrated on validation (ECE
well below 0.05), yet all three static-0.5 baselines performed nearly
identically to approve-all. The argmin-of-expected-cost policy makes the
tradeoff between fraud loss, review cost, and false-positive friction
explicit at every transaction. The mechanism is amount scaling: when the
fraud-loss term is per-row, the expected cost of approving a high-amount
transaction grows in proportion to that amount, and the argmin correctly
routes large transactions toward review even when the classifier's
probability is unchanged. Under constant loss the same transactions would
have been approved.

This is the paper's central methodological claim. A classifier is better if
and only if the policy it feeds produces lower realized cost. The results
support that claim: the same classifiers that perform near-identically to
approve-all at a 0.5 threshold produce a 57.91\% cost reduction when fed
into the argmin rule. The decision rule, not the model, does the work.

\subsection{Error Analysis and Tradeoffs}

The policy's operational recall is 0.5051 and its operational precision is
0.0654. That is, the policy flags about half of all fraud for review or
block, and roughly one in fifteen flagged transactions is actually
fraudulent. The precision figure looks low in isolation, but it is the
correct precision for a cost-sensitive rule under the frozen cost matrix:
the policy is willing to accept a low precision because the cost of a false
positive (0.1) is much smaller than the cost of a missed fraud (fraud loss
$\times$ amount). A rule that optimized precision would route fewer
transactions to review and would miss more fraud, raising realized cost.

The dominant error type at the classifier's 0.5 cut is false negatives---
2,843 of them against only 5 false positives on the test window. The policy
mitigates this not by retraining the classifier but by using its
probabilities differently: scores that would be classified as negative at
0.5 are routed to review if their expected cost justifies it. The
classifier's ranking ability, captured by its ROC-AUC of 0.8753, is
preserved; only the thresholding behavior changes.

\subsection{Connection to the Decision-Centric Framing}

The results support the decision-centric framing on three levels. First,
the classifier comparison is a non-finding: the three traditional
classifiers produce statistically indistinguishable policies, so classifier
choice does not determine performance. Second, the static-threshold
baselines perform near approve-all even though the classifiers are
well-calibrated, demonstrating that calibration alone does not produce good
decisions. Third, the argmin policy, using the same classifier
probabilities, reduces cost by 57.91\%, showing that the decision rule is
where the value is created.

%=====================================================================
\section{Limitations}
%=====================================================================

The evaluation in this paper is a controlled temporal backtest on a single
synthetic dataset under a single delay regime. It establishes the
methodological claim---that a cost-sensitive policy can substantially
reduce realized cost relative to static baselines---but it does not
establish general fraud-detection performance in production. The
limitations fall into three groups: data, assumptions, and framing.

\emph{Data limitations.} The dataset is synthetic: BAF is a realistic
benchmark, but its fraud rate, feature distributions, and dispute dynamics
are generated rather than observed, so the reported cost reductions are
not production estimates. The delay regime is a single one-month interval
determined by BAF's month-level granularity: longer delays, shorter
delays, and mixed per-transaction delays are not evaluated, and the
censored fraction is a direct consequence of the fixed regime. The amount
column is a proxy---\texttt{proposed\_credit\_limit}, not a true
transaction amount---so the amount scaling that drives much of the
policy advantage is applied to a proxy feature and its absolute magnitude
should not be taken literally.

\emph{Assumption limitations.} The cost matrix is a frozen assumption
rather than a measured quantity; the sensitivity sweep bounds how much
the result moves under 2$\times$ variation in each parameter, but the
values themselves are chosen, not estimated from operational data. The
reported classifier selection is itself an assumption: the comparison
across three classifiers was a non-finding---the top two were within
1.07\% of each other in mean cross-validated cost---so the choice of
Logistic Regression is justified by interpretability and speed, not by
superior performance. On a larger sample, a different seed, or a
different temporal fold structure, a different classifier might be
selected.

\emph{Framing limitations.} The policy does not model review capacity:
real fraud operations have a finite number of investigators per time
window, and the current build does not enforce a hard cap on the review
queue. The policy does not model feedback effects: in production, the
actions a system takes influence which labels are observed later, and
the current backtest assumes a static label set. The policy does not
adapt to drift: the training and test windows are consecutive months,
and the increase in fraud rate from training to test is documented but
not corrected.

These limitations are named as future work rather than claimed as solved.

%=====================================================================
\section{Conclusion and Future Work}
%=====================================================================

This paper evaluated a cost-sensitive fraud decision policy under a delayed
and censored label regime on the Bank Account Fraud dataset. The classifier
is an input; the policy is the object of study; the backtest is the
evidence. Under a 1-month delay, a chronological split, and a frozen cost
matrix with amount-scaled fraud loss, the argmin-of-expected-cost policy
reduced realized cost per transaction by \textbf{57.91\%} relative to the
strongest baseline, with a 95\% bootstrap confidence interval of
\textbf{[53.95\%, 62.13\%]} that excludes zero. The result is robust to a
2$\times$ variation in each of three cost parameters, with a minimum
advantage of \textbf{47.76\%}. The classifier selection across Logistic
Regression, Random Forest, and LightGBM was a non-finding under a
pre-registered noise-band guard, and the tie was broken toward the simplest
model. Censored labels (96,843 rows, 9.68\%) were excluded and reported
rather than treated as negatives.

The paper's methodological contribution is a reproducible evaluation
protocol for fraud decision policies under delayed and partial labels:
temporal splitting, a calibration gate, a pre-registered noise-band guard
for classifier selection, bootstrap confidence intervals on the primary
metric, and a cost-sensitivity sweep that bounds the robustness of the
finding. Its empirical contribution is the demonstration that static
thresholds on highly imbalanced fraud data perform near-identically to
approve-all even when the underlying classifier is well-calibrated, and
that the value of a decision system lies in the decision rule rather than
in the classifier alone. Its negative contribution is an honest reporting
of a classifier comparison that produced no finding, in place of a spurious
ranking.

Several directions remain open. The most immediate is the extension to
additional delay regimes---2-month, 3-month, and mixed per-transaction
delays---which requires a dataset with finer timestamp granularity than
BAF provides. A second direction is capacity-aware decisioning, which would
replace the unconstrained argmin with a policy that respects a finite
review budget and reports the value lost to that constraint. A third is
cost-sensitive training, in which the classifier's loss function
incorporates the same asymmetry the policy uses, evaluated against the
current decoupled design under the same stop criteria. Further directions
include rolling-window evaluation to detect drift, fairness-aware policy
constraints for regulatory contexts, and an evaluation of amount-scaled
\texttt{false\_positive\_cost}, which was named as a scope boundary in this
paper and would require a corresponding change to the cost matrix. Each is
gated on a measured failure of the current system, and none is a claim of
this submission.

%=====================================================================
% APPENDICES
%=====================================================================
\appendices

\section{Data Dictionary}

The full feature-by-feature schema---names, types, units, allowed values,
and source-time availability---is provided in
\texttt{documentation/data\_dictionary.md}. The 28 model features are 23
numeric and 5 categorical. The exclusion list is
\texttt{transaction\_id}, \texttt{fraud\_bool}, \texttt{month},
\texttt{label\_month}, \texttt{observed}, \texttt{amount\_proxy},
\texttt{proposed\_credit\_limit}, and \texttt{device\_fraud\_count}.

\section{Repository and Application}

Source code, configuration, and reproducibility instructions are available
in the project repository at
\url{https://github.com/kira-ml/delayed-label-fraud-decisioning}.

The deployed decision-support application is available at:

\url{https://delayed-label-fraud-decisioning-gefp9s9mbkfdyzhhvescdm.streamlit.app}

The application loads the same classifier and preprocessing pipeline
reported in this paper, and displays the predicted fraud probability along
with the routed action and the cost reasoning behind it.

\section{Contribution Record}

See \texttt{documentation/contribution\_record.md}.

\section{Ownership and Authorship Declaration}

See \texttt{documentation/ownership\_declaration.md}.

%=====================================================================
% REFERENCES — IEEE numbered, ordered by first in-text appearance
%=====================================================================
\begin{thebibliography}{99}

\bibitem{bsp2023roadmap}
Bangko Sentral ng Pilipinas, \emph{Digital Payments Transformation
Roadmap 2020--2023}, Manila, Philippines, 2023.

\bibitem{bspspei}
Bangko Sentral ng Pilipinas, \emph{Selected Philippine Economic
Indicators --- Payments}, Manila, Philippines, 2024.

\bibitem{bspcam}
Bangko Sentral ng Pilipinas, \emph{Consumer Assistance Mechanism Annual
Report}, Manila, Philippines, 2023.

\bibitem{visa}
Visa Inc., \emph{Dispute Resolution Rules}, public materials, 2023.

\bibitem{chapelle2014}
O. Chapelle, ``Modeling delayed feedback in display advertising,'' in
\emph{Proc. 20th ACM SIGKDD Int. Conf. Knowledge Discovery and Data
Mining}, 2014, pp. 1097--1105.

\bibitem{elkan2001}
C. Elkan, ``The foundations of cost-sensitive learning,'' in \emph{Proc.
17th Int. Joint Conf. Artificial Intelligence}, 2001, pp. 973--978.

\bibitem{jesus2022}
S. Jesus, J. Pombal, D. Alves, A. Cruz, P. Saleiro, R. Ribeiro, J. Gama,
and P. Bizarro, ``Turning the tables: Biased, imbalanced, dynamic tabular
datasets for ML evaluation,'' in \emph{Proc. NeurIPS Datasets and
Benchmarks Track}, 2022.

\bibitem{elkan2008}
C. Elkan and K. Noto, ``Learning classifiers from only positive and
unlabeled data,'' in \emph{Proc. 14th ACM SIGKDD Int. Conf. Knowledge
Discovery and Data Mining}, 2008, pp. 213--220.

\bibitem{dalpozzolo2015}
A. Dal Pozzolo, O. Caelen, R. A. Johnson, and G. Bontempi, ``Calibrating
probability with an unbalanced class: An application to fraud detection,''
in \emph{Proc. IEEE Symp. Computational Intelligence and Data Mining},
2015, pp. 1--7.

\bibitem{pedregosa2011}
F. Pedregosa \emph{et al.}, ``Scikit-learn: Machine learning in Python,''
\emph{J. Mach. Learn. Res.}, vol. 12, pp. 2825--2830, 2011.

\bibitem{ke2017}
G. Ke \emph{et al.}, ``LightGBM: A highly efficient gradient boosting
decision tree,'' in \emph{Advances in Neural Information Processing
Systems}, vol. 30, 2017, pp. 3146--3154.

\end{thebibliography}

\end{document}